% PAIRMATCH3-REVIEW
\documentclass[11pt]{article}
\usepackage[margin=.88in]{geometry}
\usepackage{amsmath,amssymb,amsthm}
\begin{document}
\begin{center}
{\Large PAIRMATCH3-REVIEW: literature check after pair-matched order three}\\[4pt]
\textbf{Date: 2026-10-04 20:34 PT}\\
\textbf{Status: REVIEW, no new computation}\\
\textbf{P versus NP status: OPEN}
\end{center}

\section{Scope}
This note records an external review of
\texttt{2026-10-04\_2250\_PAIRMATCH3.tex}, read against the concurrent exact census in
\texttt{2026-10-04\_2355\_ORDER3EXACT1.tex}.
No new formulas were sampled. No asymptotic claim is added.

\section{What the runs support}
After nearest-neighbor matching on the six dominant pair coordinates, held-out pair balanced accuracy sits in \(0.50\)--\(0.58\). Under the caliper \(d_{\mathrm{pair}}\le 0.75\), residual mismatch is about \(0.09\)--\(0.24\) SAT standard deviations per coordinate, and the sampled triple-orbit scan is
\[
0.548,\ 0.500,\ 0.700,\ 0.571
\]
at \(q=20,24,28,32\), on \(41,32,29,13\) retained positives. The isolated \(0.70\) at \(q=28\) is compatible with noise and does not replicate at neighboring sizes.

The exact census does not repair that. At \(q=20\), \(t=70\), \(N=3000\), twenty caliper survivors, all \(\binom{70}{3}=54740\) triples, \(81\) canonical orbits, held-out balanced accuracy \(0.45\) on ten evaluation pairs. The \(q=16\) score \(0.643\) is not reproduced.

\[
\boxed{\text{no replicated order-three distinguisher after pair matching}}
\]
is the correct computational box. It is weaker than order-three indistinguishability, and it is not a monotone-circuit lower bound.

Deletion minima are not MUS certificates. The stronger \(q=20\) pass found \(\overline m_{\mathrm{core}}=31.30\), \(m_{\min}=26\), \(\overline n_{\mathrm{core}}=17.85\), \(n_{\min}=15\), over \(50\) random deletion orders.

\section{Literature alignment}
The working density \(t=3.5q\) lies above the pure-literal threshold and below the satisfiability threshold. Scott--Sorkin prove that below the pure-literal threshold a rare unsatisfiable \(r\)-SAT formula has, with high probability, a unique minimal unsatisfiable subformula of minimum excess, and only finitely many such cores of each excess. Above the satisfiability threshold, Chv\'atal--Szemer\'edi says the obstruction is linear. Their Question~1 --- whether large minimal unsatisfiable subformulas remain unlikely for \(\alpha\le(1-\varepsilon)\alpha_{\mathrm{sat}}\) --- is exactly the open interval containing \(\alpha=3.5\). The random \(3\)-SAT pure-literal point is about \(1.636\), as already used in SLICEDENSITY2. A bulk three-clause orbit census is the wrong aggregator for a sparse core, and sampling \(400\)--\(600\) triples out of \(\binom{t}{3}\) makes that worse.

The positive SAT/UNSAT classifiers in the literature are not evidence against this null. Xu--Hoos--Leyton-Brown obtain about \(70\%\) at the phase transition, \(\alpha\approx 4.26\), and their two-feature tree is LPSLACK coefficient of variation plus positive/negative balance: an LP-margin feature and a polarity feature, not a triple-orbit count. Cameron--Chen--Hartford--Leyton-Brown reach roughly \(79\)--\(84\%\) on \(100\)--\(600\) variables with exchangeable and message-passing nets, still at the transition, still without exhibiting a local motif. Zhou's long-range frustration order appears only for \(\alpha>\alpha_R(3)\approx 4.19\); whitening separates local frozen cores from frustration loops of length \(\log n\). Neither is the expected bulk reason for rare UNSAT at \(\alpha=3.5\). Bresler--Huang's many-way overlap gap is an algorithmic barrier for finding solutions at large \(k\), far above this density. Hypergraph motifs on three connected hyperedges are the unsigned cousin of TRIPLE-ORBIT-CANON; they separate domains of real hypergraphs from null models, not SAT from UNSAT.

Sunflower machinery is the wrong obstruction if the core is small. Razborov's approximation method, and the robust-sunflower strengthenings of Alweiss--Lovett--Wu--Zhang and Cavalar--Kumar--Rossman, lower-bound monotone circuits by plucking large uniform set systems. A minimum-excess MUF is the opposite object. SAT-SUNFLOWER2 should wait on the core-size measurement. Filmus--Mossel slice invariance and the Filmus--Ihringer junta theorem, already cited in ORDER3EXACT1, make the degree hierarchy meaningful and still do not imply that a small general monotone circuit has low harmonic degree.

\section{Resource that still has to scale}
Grant unlimited polynomial-dimensional attention, unlimited random access, exact Gaussian elimination, exact \(2\)-SAT, and polynomial-depth global attention. Pair matching prices in what exact \(2\)-SAT and the implication graph see. The orbit scan prices in depth-\(1\) attention on every signed triple. Both went quiet. Two candidates remain.

\begin{enumerate}
\item A small excess core. The incompatibility is local, but which variables participate is not a low-order marginal. Attention represents the search for the subsystem; it does not escape it. If excess stays bounded, certifying that no small core exists need not grow with \(q\).
\item A delocalized incompatibility on the matched pair. The symmetric difference of the matched SAT and UNSAT formulas is spread over \(\omega(1)\) variables, and LP slack --- the Xu feature, already free under exact Gaussian elimination --- also fails to separate. That is the synchronization bottleneck: compatibility of overlapping local charts, not another chart.
\end{enumerate}
The caliper table cannot tell these apart.

\section{Revised queue}
\begin{enumerate}
\item \textbf{MATCHED-MUF2 first}, on pairs already in hand. Exact or branch-and-bound MUS at \(q=16,18,20\). Record excess \(e=2m-n_{\mathrm{core}}\), uniqueness, and whether the core is absent from the matched SAT twin. Scott--Sorkin say minimum excess is the right coordinate, not raw clause count.
\item \textbf{Free-resource check} before any new generator. On the same matched pairs, compute LPSLACK coefficient of variation and the GE/XOR defect. If slack separates, the residual is inside resources already granted.
\item \textbf{ORDER3EXACT2 only inside the core}, or inside the failed-literal closure, plus one confirmatory bulk pass at \(q=20\). A null bulk census is the prediction if MATCHED-MUF2 finds a small core. Do not lead with \(q=32\) bulk enumeration.
\item \textbf{PAIR-MATCH2} only after that, and only at \(q=28\), with on the order of a hundred caliper-retained positives, to confirm or kill the \(0.70\).
\item \textbf{Hold SAT-SUNFLOWER2} and order-four statistics until the core-order distribution is in hand. A linear, delocalized core at \(\alpha=3.5\) would be a structural fact inside Scott--Sorkin's open interval; only then is a higher-order or sunflower statistic aimed at the resource that still scales.
\end{enumerate}

\section{Status}
TRIPLE-ORBIT-CANON and MOMENT-CAMOUFLAGE remain candidate lemmas with elementary proofs, as recorded in PAIRMATCH3 and ORDER3EXACT1. This review is not a computation and adds no theorem. Asymptotic degree-three camouflage, a certified minimum-core lower bound, and the superpolynomial monotone slice lower bound stay \textbf{OPEN}. Therefore \(P\ne NP\) remains OPEN.

\begin{thebibliography}{9}
\bibitem{SS2010} A.~D. Scott and G.~B. Sorkin, Structure of random \(r\)-SAT below the pure literal threshold, arXiv:1008.1260.
\bibitem{CS1988} V.~Chv\'atal and E.~Szemer\'edi, Many hard examples for resolution, J.~ACM 35 (1988).
\bibitem{XHL2012} L.~Xu, H.~H. Hoos, and K.~Leyton-Brown, Predicting satisfiability at the phase transition, AAAI 2012.
\bibitem{CCHL2019} C.~Cameron, R.~Chen, J.~Hartford, and K.~Leyton-Brown, Predicting propositional satisfiability via end-to-end learning, AAAI 2019.
\bibitem{Z2005} H.~Zhou, Long-range frustration in finite connectivity spin glasses and random \(K\)-satisfiability, New J.~Phys.\ 7 (2005) 123.
\bibitem{Z2009} H.-J. Zhou, From one solution of a 3-satisfiability formula to a solution cluster, Phys.~Rev.~E 79 (2009) 031102.
\bibitem{BH2021} G.~Bresler and B.~Huang, The algorithmic phase transition of random \(k\)-SAT for low-degree polynomials, FOCS 2021.
\bibitem{ALWZ2021} R.~Alweiss, S.~Lovett, K.~Wu, and J.~Zhang, Improved bounds for the sunflower lemma, Ann.\ of Math.\ 194 (2021).
\bibitem{CKR2022} B.~P. Cavalar, M.~Kumar, and B.~Rossman, Monotone circuit lower bounds from robust sunflowers, Algorithmica 2022.
\end{thebibliography}

\end{document}
